%% -*- coding:utf-8 -*-

\chapter{Ausblick}
\label{chap-outlook}

In diesem Buch wurden Fragmente von Grammatiken mehrerer \ili{germanisch}er Sprachen skizziert. Die im
Hintergrund stehende Theorie ist die Head-Driven Phrase Structure Grammar (HPSG)\indexhpsg
\parencites{ps}{ps2}{MuellerLehrbuch3}{HPSGHandbook}. Wir haben uns die Valenz angesehen und wie sie
in Valenzlisten wie \spr und \comps repräsentiert wird. Wir haben uns auch Adjunkte angesehen, die nicht
in Listen repräsentiert werden: Adjunkte selegieren die Köpfe, die sie modifizieren. HPSG nimmt an, dass es Schemata
für die Kombination von sprachlichem Material gibt. Wir haben uns mit dem Spezifikator-Kopf-Schema, dem
Kopf-Komplement-Schema, dem Kopf-Adjunkt-Schema und auch mit dem Prädikatskomplex-Schema befasst. Verbalkomplexe
in den \ili{germanisch}en OV-Sprachen wurden als Prädikatskomplexbildung analysiert.

Die \ili{germanisch}en Sprachen variieren hinsichtlich ihrer Grundabfolge (VO oder OV). Mit Ausnahme des \ili{Englisch}en
sind alle \ili{germanisch}en Sprachen V2-Sprachen. V2-Sätze werden über Kopfbewegung analysiert: Es gibt ein
leeres Verb in Endstellung, das mit dem vorangestellten Verb in Beziehung steht.

Alle Analysen sind in computerverarbeitbaren Grammatikfragmenten implementiert. Sie sind vollständig formalisiert –
andernfalls wären sie nicht verarbeitbar –, wurden hier jedoch in vereinfachter und skizzenhafter
Form wiedergegeben. Ich habe in Abschnitt~\ref{sec-linking} kurz über den Zusammenhang zwischen Syntax und Semantik gesprochen, aber
natürlich gehen alle Implementierungen mit semantischen Repräsentationen einher.

Aus Platzgründen ist es nicht möglich, alle Konzepte der HPSG sorgfältig einzuführen, aber die
interessierte Leserin sei eingeladen, einen Blick in die HPSG-Monographien
\parencites{ps}{ps2}{GSag2000a-u}{MuellerLehrbuch3}, in die Überblicksartikel
\parencites{LM2006a}{PK2006a-u}{Bildhauer2014a-u}{MuellerHPSGHandbook,MuellerCurrentApproaches}, in Kapitel~9 des
Lehrbuchs zur Grammatiktheorie \citep{MuellerGT-Eng5} oder in das Handbuch zur HPSG
\citep{HPSGHandbook} zu werfen. Besonders der letztgenannte Band ist ein aktuelles Buch mit mehr als 1600 Seiten,
das sich in 32 Kapiteln mit fast jedem Aspekt befasst, für den man sich interessieren könnte.







%      <!-- Local IspellDict: de_DE -->
